% Lösungen Einheit 3 von 3 – Waagerechte Asymptoten (zusammenbau v0.1)
\einheitenkopf{Einheit 3 von 3 · Waagerechte Asymptoten}

\erg{15}{a) der Summand $-7$ \quad b) $f(x) \to 4$ für $x \to +\infty$; waagerechte Asymptote $y = 4$ \quad c) Grenzwert $-4$; streng monoton fallend wie $e^{-0{,}5x}$, weil der Streckfaktor positiv ist und die Verschiebung die Monotonie nicht ändert; $f(0) = 0$, also durch den Ursprung. \quad d) Grenzwert $8$ für $x \to -\infty$ (Nenner gegen eins), Grenzwert null für $x \to +\infty$ (Nenner wächst); keine Nullstelle, weil der Zähler nie null ist. \quad e) Grenzwert $4$ für $t \to +\infty$: die Kühlschranktemperatur, der sich das Getränk auf Dauer nähert; $T$ fällt streng monoton mit $T(0) = 22$, Wertebereich $4 < T(t) \le 22$.}
\erg{16}{a) Fehler: die Verschiebung $-9$ vergessen. Richtig: $f(x) \to -9$, Asymptote $y = -9$. \quad b) $e^{-x}$ wird beliebig klein, also auch $5e^{-x}$; übrig bleibt der Summand $3$. \quad c) $T(t) \to 20$: die Raumtemperatur; kälter als der Raum wird der Kuchen nicht.}
